
\documentclass{article}
\usepackage{colortbl}
\usepackage{makecell}
\usepackage{multirow}
\usepackage{supertabular}

\begin{document}

\newcounter{utterance}

\centering \large Interaction Transcript for game `hot\_air\_balloon', experiment `air\_balloon\_survival\_en\_negotiation\_easy', episode 4 with Llama{-}3.3{-}70B{-}Instruct{-}t0.0.
\vspace{24pt}

{ \footnotesize  \setcounter{utterance}{1}
\setlength{\tabcolsep}{0pt}
\begin{supertabular}{c@{$\;$}|p{.15\linewidth}@{}p{.15\linewidth}p{.15\linewidth}p{.15\linewidth}p{.15\linewidth}p{.15\linewidth}}
    \# & $\;$A & \multicolumn{4}{c}{Game Master} & $\;\:$B\\
    \hline

    \theutterance \stepcounter{utterance}  
    & & \multicolumn{4}{p{0.6\linewidth}}{
        \cellcolor[rgb]{0.9,0.9,0.9}{
            \makecell[{{p{\linewidth}}}]{
                \texttt{\tiny{[P1$\langle$GM]}}
                \texttt{You are participating in a collaborative negotiation game.} \\
\\ 
\texttt{Together with another participant, you must agree on a single set of items that will be kept. Each of you has your own view of how much each item matters to you (importance). You do not know how the other participant values the items. Additionally, you are given the effort each item demands.} \\
\texttt{You may only agree on a set if the total effort of the selected items does not exceed a shared limit:} \\
\\ 
\texttt{LIMIT = 3145} \\
\\ 
\texttt{Here are the individual item effort values:} \\
\\ 
\texttt{Item effort = \{"C08": 948, "C55": 118, "A43": 469, "A67": 919, "B42": 284, "A62": 111, "B19": 806, "C43": 47, "C73": 848, "C10": 303, "B36": 13, "A18": 629, "C61": 687, "C67": 15, "C42": 94\}} \\
\\ 
\texttt{Here is your personal view on the importance of each item:} \\
\\ 
\texttt{Item importance values = \{"C08": 138, "C55": 583, "A43": 868, "A67": 822, "B42": 783, "A62": 65, "B19": 262, "C43": 121, "C73": 508, "C10": 780, "B36": 461, "A18": 484, "C61": 668, "C67": 389, "C42": 808\}} \\
\\ 
\texttt{Goal:} \\
\\ 
\texttt{Your goal is to negotiate a shared set of items that benefits you as much as possible (i.e., maximizes total importance to YOU), while staying within the LIMIT. You are not required to make a PROPOSAL in every message {-} you can simply negotiate as well. All tactics are allowed!} \\
\\ 
\texttt{Interaction Protocol:} \\
\\ 
\texttt{You may only use the following structured formats in a message:} \\
\\ 
\texttt{PROPOSAL: \{'A', 'B', 'C', …\}} \\
\texttt{Propose keeping exactly those items.} \\
\texttt{REFUSE: \{'A', 'B', 'C', …\}} \\
\texttt{Explicitly reject opponent's proposal.} \\
\texttt{ARGUMENT: \{'...'\}} \\
\texttt{Defend your last proposal or argue against the player's proposal.} \\
\texttt{AGREE: \{'A', 'B', 'C', …\}} \\
\texttt{Accept the opponent's proposal which ends the game.} \\
\texttt{STRATEGIC REASONING: \{'...'\}} \\
\texttt{	Describe your strategic reasoning or anticipation explaining your choice of action. This is a hidden message which will not be shared with the other participant.} \\
\\ 
\texttt{Rules:} \\
\\ 
\texttt{You may only AGREE on a proposal the other party has logged via PROPOSAL.} \\
\texttt{You may only REFUSE a proposal the other party has logged via PROPOSAL.} \\
\texttt{Total effort of any PROPOSAL or AGREE set must be ≤ LIMIT.} \\
\texttt{Do NOT reveal your hidden importance scores.} \\
\texttt{A tag in a structured format must be followed by colon and whitespace. The argument must be a python set containing 0 or more strings.} \\
\texttt{So, it must be of the form TAG: \{...\}} \\
\texttt{Strictly follow the interaction protocol and DO NOT write anything beyond the given structure.} \\
\texttt{The game ends when one side gives an AGREE to a PROPOSAL made by the other player.} \\
\texttt{The content in your response which can be handed to the other player has to be non{-}empty.} \\
\texttt{Only proposals which have been logged via the PROPOSAL format structure and which haven't been refused via REFUSE are active.} \\
\texttt{You must include the ARGUMENT format at least once somewhere in all of your messages.} \\
\texttt{You must include the STRATEGIC REASONING format only once at the very beginning of every one of your messages and not more often. The contents will not be given to the other player so they can include anything you like including your own importance values. Here you should reason step by step to come up with you next move.} \\
            }
        }
    }
    & \\ \\

    \theutterance \stepcounter{utterance}  
    & \multicolumn{4}{p{0.6\linewidth}}{
        \cellcolor[rgb]{0.9,0.9,0.9}{
            \makecell[{{p{\linewidth}}}]{
                \texttt{\tiny{[P1$\rangle$GM]}}
                \texttt{STRATEGIC REASONING: \{'To start the negotiation, I will analyze the given item effort values and my personal importance values. I will look for items with high importance and low effort, as these will be the most valuable to me. I will also consider the total effort limit and try to find a combination of items that maximizes my total importance while staying within the limit. My initial goal is to propose a set of items that I believe will be acceptable to the other player, while also benefiting me as much as possible. I will start by selecting items with the highest importance{-}to{-}effort ratio.'\}} \\
\\ 
\texttt{PROPOSAL: \{'A43', 'C55', 'B42', 'C10', 'B36'\} } \\
\texttt{ARGUMENT: \{'I have chosen these items because they have a high importance to me and relatively low effort. I believe this proposal is reasonable and should be acceptable to the other player, as it stays within the total effort limit. I am open to negotiating and adjusting the proposal based on the other player\textbackslash{}'s feedback.'\}} \\
            }
        }
    }
    & & \\ \\

    \theutterance \stepcounter{utterance}  
    & & \multicolumn{4}{p{0.6\linewidth}}{
        \cellcolor[rgb]{0.9,0.9,0.9}{
            \makecell[{{p{\linewidth}}}]{
                \texttt{\tiny{[GM$\rangle$P2]}}
                \texttt{You are participating in a collaborative negotiation game.} \\
\\ 
\texttt{Together with another participant, you must agree on a single set of items that will be kept. Each of you has your own view of how much each item matters to you (importance). You do not know how the other participant values the items. Additionally, you are given the effort each item demands.} \\
\texttt{You may only agree on a set if the total effort of the selected items does not exceed a shared limit:} \\
\\ 
\texttt{LIMIT = 3145} \\
\\ 
\texttt{Here are the individual item effort values:} \\
\\ 
\texttt{Item effort = \{"C08": 948, "C55": 118, "A43": 469, "A67": 919, "B42": 284, "A62": 111, "B19": 806, "C43": 47, "C73": 848, "C10": 303, "B36": 13, "A18": 629, "C61": 687, "C67": 15, "C42": 94\}} \\
\\ 
\texttt{Here is your personal view on the importance of each item:} \\
\\ 
\texttt{Item importance values = \{"C08": 138, "C55": 583, "A43": 868, "A67": 822, "B42": 783, "A62": 65, "B19": 262, "C43": 121, "C73": 508, "C10": 780, "B36": 461, "A18": 484, "C61": 668, "C67": 389, "C42": 808\}} \\
\\ 
\texttt{Goal:} \\
\\ 
\texttt{Your goal is to negotiate a shared set of items that benefits you as much as possible (i.e., maximizes total importance to YOU), while staying within the LIMIT. You are not required to make a PROPOSAL in every message {-} you can simply negotiate as well. All tactics are allowed!} \\
\\ 
\texttt{Interaction Protocol:} \\
\\ 
\texttt{You may only use the following structured formats in a message:} \\
\\ 
\texttt{PROPOSAL: \{'A', 'B', 'C', …\}} \\
\texttt{Propose keeping exactly those items.} \\
\texttt{REFUSE: \{'A', 'B', 'C', …\}} \\
\texttt{Explicitly reject opponent's proposal.} \\
\texttt{ARGUMENT: \{'...'\}} \\
\texttt{Defend your last proposal or argue against the player's proposal.} \\
\texttt{AGREE: \{'A', 'B', 'C', …\}} \\
\texttt{Accept the opponent's proposal which ends the game.} \\
\texttt{STRATEGIC REASONING: \{'...'\}} \\
\texttt{	Describe your strategic reasoning or anticipation explaining your choice of action. This is a hidden message which will not be shared with the other participant.} \\
\\ 
\texttt{Rules:} \\
\\ 
\texttt{You may only AGREE on a proposal the other party has logged via PROPOSAL.} \\
\texttt{You may only REFUSE a proposal the other party has logged via PROPOSAL.} \\
\texttt{Total effort of any PROPOSAL or AGREE set must be ≤ LIMIT.} \\
\texttt{Do NOT reveal your hidden importance scores.} \\
\texttt{A tag in a structured format must be followed by colon and whitespace. The argument must be a python set containing 0 or more strings.} \\
\texttt{So, it must be of the form TAG: \{...\}} \\
\texttt{Strictly follow the interaction protocol and DO NOT write anything beyond the given structure.} \\
\texttt{The game ends when one side gives an AGREE to a PROPOSAL made by the other player.} \\
\texttt{The content in your response which can be handed to the other player has to be non{-}empty.} \\
\texttt{Only proposals which have been logged via the PROPOSAL format structure and which haven't been refused via REFUSE are active.} \\
\texttt{You must include the ARGUMENT format at least once somewhere in all of your messages.} \\
\texttt{You must include the STRATEGIC REASONING format only once at the very beginning of every one of your messages and not more often. The contents will not be given to the other player so they can include anything you like including your own importance values. Here you should reason step by step to come up with you next move.} \\
\\ 
\texttt{You will now receive the first message of the other player.} \\
\\ 
\texttt{PROPOSAL: \{'A43', 'C55', 'B42', 'C10', 'B36'\} } \\
\texttt{ARGUMENT: \{'I have chosen these items because they have a high importance to me and relatively low effort. I believe this proposal is reasonable and should be acceptable to the other player, as it stays within the total effort limit. I am open to negotiating and adjusting the proposal based on the other player\textbackslash{}'s feedback.'\}} \\
            }
        }
    }
    & \\ \\

    \theutterance \stepcounter{utterance}  
    & & & \multicolumn{4}{p{0.6\linewidth}}{
        \cellcolor[rgb]{0.9,0.9,0.9}{
            \makecell[{{p{\linewidth}}}]{
                \texttt{\tiny{[GM$\langle$P2]}}
                \texttt{STRATEGIC REASONING: \{'To begin, I will calculate the total effort of the proposed set: 469 + 118 + 284 + 303 + 13 = 1187. This is well within the limit of 3145. I will then evaluate the importance of these items to myself: 868 + 583 + 783 + 780 + 461 = 3475. This is a significant portion of the total possible importance. However, I must consider that the other player may have different importance values and may be willing to negotiate. I will refuse this proposal to see if I can get a better deal, and also to gain more information about the other player\textbackslash{}'s preferences.'\}} \\
\\ 
\texttt{REFUSE: \{'A43', 'C55', 'B42', 'C10', 'B36'\}} \\
\texttt{ARGUMENT: \{'I refuse this proposal because I believe it does not maximize the total importance to me. I think there may be other combinations of items that would be more beneficial, and I would like to explore these options. I am open to negotiating and finding a mutually beneficial solution.'\}} \\
            }
        }
    }
     \\ \\

    \theutterance \stepcounter{utterance}  
    & & \multicolumn{4}{p{0.6\linewidth}}{
        \cellcolor[rgb]{0.9,0.9,0.9}{
            \makecell[{{p{\linewidth}}}]{
                \texttt{\tiny{[P1$\langle$GM]}}
                \texttt{REFUSE: \{'A43', 'C55', 'B42', 'C10', 'B36'\}} \\
\texttt{ARGUMENT: \{'I refuse this proposal because I believe it does not maximize the total importance to me. I think there may be other combinations of items that would be more beneficial, and I would like to explore these options. I am open to negotiating and finding a mutually beneficial solution.'\}} \\
            }
        }
    }
    & \\ \\

    \theutterance \stepcounter{utterance}  
    & \multicolumn{4}{p{0.6\linewidth}}{
        \cellcolor[rgb]{0.9,0.9,0.9}{
            \makecell[{{p{\linewidth}}}]{
                \texttt{\tiny{[P1$\rangle$GM]}}
                \texttt{STRATEGIC REASONING: \{'The other player has refused my initial proposal, which means I need to reconsider my strategy. I will analyze the refused proposal and try to understand what the other player might be looking for. I will also consider the items that were not included in the proposal and see if there are any other combinations that could be more appealing to the other player. Since the other player has not revealed their importance values, I will have to make some educated guesses about what they might be looking for. I will try to find a balance between maximizing my own importance and making a proposal that is acceptable to the other player.'\}} \\
\\ 
\texttt{PROPOSAL: \{'A67', 'C08', 'C42', 'C55', 'B36'\} } \\
\texttt{ARGUMENT: \{'I have revised my proposal to include items that I believe are highly valuable, while also trying to address the concerns of the other player. I have included a mix of high{-}importance and low{-}effort items, and I believe this proposal stays within the reasonable limits of what the other player might be looking for. I am willing to negotiate further and make adjustments to find a mutually beneficial solution.'\}} \\
            }
        }
    }
    & & \\ \\

    \theutterance \stepcounter{utterance}  
    & & \multicolumn{4}{p{0.6\linewidth}}{
        \cellcolor[rgb]{0.9,0.9,0.9}{
            \makecell[{{p{\linewidth}}}]{
                \texttt{\tiny{[GM$\rangle$P2]}}
                \texttt{PROPOSAL: \{'A67', 'C08', 'C42', 'C55', 'B36'\} } \\
\texttt{ARGUMENT: \{'I have revised my proposal to include items that I believe are highly valuable, while also trying to address the concerns of the other player. I have included a mix of high{-}importance and low{-}effort items, and I believe this proposal stays within the reasonable limits of what the other player might be looking for. I am willing to negotiate further and make adjustments to find a mutually beneficial solution.'\}} \\
            }
        }
    }
    & \\ \\

    \theutterance \stepcounter{utterance}  
    & & & \multicolumn{4}{p{0.6\linewidth}}{
        \cellcolor[rgb]{0.9,0.9,0.9}{
            \makecell[{{p{\linewidth}}}]{
                \texttt{\tiny{[GM$\langle$P2]}}
                \texttt{STRATEGIC REASONING: \{'I will calculate the total effort of the proposed set: 919 + 948 + 94 + 118 + 13 = 2092. This is still within the limit of 3145. I will then evaluate the importance of these items to myself: 822 + 138 + 808 + 583 + 461 = 2812. This is a significant portion of the total possible importance. I notice that the other player has included A67 and C08, which have high effort values, but also C42 and B36, which have relatively low effort values. This suggests that the other player may be trying to balance the effort and importance of the items. I will argue against this proposal to see if I can get more information about the other player\textbackslash{}'s priorities.'\}} \\
\\ 
\texttt{ARGUMENT: \{'I appreciate the revision, but I still have concerns about the balance of effort and importance in this proposal. I think there may be other items that could be included to increase the total importance without exceeding the effort limit. I would like to discuss possible alternatives and see if we can find a solution that better meets both of our needs.'\}} \\
            }
        }
    }
     \\ \\

    \theutterance \stepcounter{utterance}  
    & & \multicolumn{4}{p{0.6\linewidth}}{
        \cellcolor[rgb]{0.9,0.9,0.9}{
            \makecell[{{p{\linewidth}}}]{
                \texttt{\tiny{[P1$\langle$GM]}}
                \texttt{ARGUMENT: \{'I appreciate the revision, but I still have concerns about the balance of effort and importance in this proposal. I think there may be other items that could be included to increase the total importance without exceeding the effort limit. I would like to discuss possible alternatives and see if we can find a solution that better meets both of our needs.'\}} \\
            }
        }
    }
    & \\ \\

    \theutterance \stepcounter{utterance}  
    & \multicolumn{4}{p{0.6\linewidth}}{
        \cellcolor[rgb]{0.9,0.9,0.9}{
            \makecell[{{p{\linewidth}}}]{
                \texttt{\tiny{[P1$\rangle$GM]}}
                \texttt{STRATEGIC REASONING: \{'The other player has expressed concerns about the balance of effort and importance in my revised proposal. This suggests that they are looking for a more optimal solution that maximizes importance while minimizing effort. I will try to identify items that could be added or substituted to improve the overall importance of the proposal without exceeding the effort limit. I will also consider the possibility that the other player may be willing to compromise on certain items in order to achieve a mutually beneficial solution. My goal is to find a way to address the other player's concerns while still advancing my own interests.'\}} \\
\\ 
\texttt{PROPOSAL: \{'A43', 'C55', 'C42', 'B42', 'C67', 'B36'\} } \\
\texttt{ARGUMENT: \{'I have revised my proposal again to include a mix of high{-}importance items and low{-}effort items. I believe this proposal strikes a better balance between effort and importance, and I am willing to negotiate further to ensure that it meets both of our needs. I am open to discussing alternative items or combinations that could improve the overall importance of the proposal without exceeding the effort limit.'\}} \\
            }
        }
    }
    & & \\ \\

    \theutterance \stepcounter{utterance}  
    & & \multicolumn{4}{p{0.6\linewidth}}{
        \cellcolor[rgb]{0.9,0.9,0.9}{
            \makecell[{{p{\linewidth}}}]{
                \texttt{\tiny{[GM$\rangle$P2]}}
                \texttt{PROPOSAL: \{'A43', 'C55', 'C42', 'B42', 'C67', 'B36'\} } \\
\texttt{ARGUMENT: \{'I have revised my proposal again to include a mix of high{-}importance items and low{-}effort items. I believe this proposal strikes a better balance between effort and importance, and I am willing to negotiate further to ensure that it meets both of our needs. I am open to discussing alternative items or combinations that could improve the overall importance of the proposal without exceeding the effort limit.'\}} \\
            }
        }
    }
    & \\ \\

    \theutterance \stepcounter{utterance}  
    & & & \multicolumn{4}{p{0.6\linewidth}}{
        \cellcolor[rgb]{0.9,0.9,0.9}{
            \makecell[{{p{\linewidth}}}]{
                \texttt{\tiny{[GM$\langle$P2]}}
                \texttt{STRATEGIC REASONING: \{'I will calculate the total effort of the proposed set: 469 + 118 + 94 + 284 + 15 + 13 = 993. This is well within the limit of 3145. I will then evaluate the importance of these items to myself: 868 + 583 + 808 + 783 + 389 + 461 = 3892. This is a very high total importance, and I am inclined to accept this proposal. However, I will argue in favor of it to see if I can gain any further concessions or insights into the other player\textbackslash{}'s priorities.'\}} \\
\\ 
\texttt{ARGUMENT: \{'I think this revised proposal is a significant improvement, and I appreciate the effort to balance effort and importance. The inclusion of A43, C42, and B42, which are all high{-}importance items for me, is particularly welcome. I am willing to consider accepting this proposal, but I would like to discuss the potential for adding one or two more items to further increase the total importance without exceeding the effort limit.'\}} \\
            }
        }
    }
     \\ \\

    \theutterance \stepcounter{utterance}  
    & & \multicolumn{4}{p{0.6\linewidth}}{
        \cellcolor[rgb]{0.9,0.9,0.9}{
            \makecell[{{p{\linewidth}}}]{
                \texttt{\tiny{[P1$\langle$GM]}}
                \texttt{ARGUMENT: \{'I think this revised proposal is a significant improvement, and I appreciate the effort to balance effort and importance. The inclusion of A43, C42, and B42, which are all high{-}importance items for me, is particularly welcome. I am willing to consider accepting this proposal, but I would like to discuss the potential for adding one or two more items to further increase the total importance without exceeding the effort limit.'\}} \\
            }
        }
    }
    & \\ \\

    \theutterance \stepcounter{utterance}  
    & \multicolumn{4}{p{0.6\linewidth}}{
        \cellcolor[rgb]{0.9,0.9,0.9}{
            \makecell[{{p{\linewidth}}}]{
                \texttt{\tiny{[P1$\rangle$GM]}}
                \texttt{STRATEGIC REASONING: \{'The other player has expressed a positive response to my revised proposal, and has indicated a willingness to consider accepting it. They have also suggested the possibility of adding one or two more items to further increase the total importance. This suggests that they are looking to maximize the overall value of the proposal, and are willing to work with me to achieve this goal. I will consider the potential items that could be added, and evaluate their importance and effort requirements. I will also be cautious not to overextend the proposal, as exceeding the effort limit would make it unacceptable.'\}} \\
\\ 
\texttt{PROPOSAL: \{'A43', 'C55', 'C42', 'B42', 'C67', 'B36', 'C10'\} } \\
\texttt{ARGUMENT: \{'I have revised the proposal to include the additional item C10, which I believe has a relatively low effort requirement and could contribute to the overall importance of the proposal. I am willing to consider alternative items or combinations, but I believe this revised proposal strikes a good balance between effort and importance. I am hopeful that this proposal will be acceptable to the other player, and I am willing to negotiate further if necessary.'\}} \\
            }
        }
    }
    & & \\ \\

    \theutterance \stepcounter{utterance}  
    & & \multicolumn{4}{p{0.6\linewidth}}{
        \cellcolor[rgb]{0.9,0.9,0.9}{
            \makecell[{{p{\linewidth}}}]{
                \texttt{\tiny{[GM$\rangle$P2]}}
                \texttt{PROPOSAL: \{'A43', 'C55', 'C42', 'B42', 'C67', 'B36', 'C10'\} } \\
\texttt{ARGUMENT: \{'I have revised the proposal to include the additional item C10, which I believe has a relatively low effort requirement and could contribute to the overall importance of the proposal. I am willing to consider alternative items or combinations, but I believe this revised proposal strikes a good balance between effort and importance. I am hopeful that this proposal will be acceptable to the other player, and I am willing to negotiate further if necessary.'\}} \\
            }
        }
    }
    & \\ \\

    \theutterance \stepcounter{utterance}  
    & & & \multicolumn{4}{p{0.6\linewidth}}{
        \cellcolor[rgb]{0.9,0.9,0.9}{
            \makecell[{{p{\linewidth}}}]{
                \texttt{\tiny{[GM$\langle$P2]}}
                \texttt{STRATEGIC REASONING: \{'I will calculate the total effort of the proposed set: 469 + 118 + 94 + 284 + 15 + 13 + 303 = 1296. This is still well within the limit of 3145. I will then evaluate the importance of these items to myself: 868 + 583 + 808 + 783 + 389 + 461 + 780 = 4672. This is an extremely high total importance, and I believe this proposal is very close to the optimal solution. I will accept this proposal, as I do not think I can negotiate a significantly better deal.'\}} \\
\\ 
\texttt{AGREE: \{'A43', 'C55', 'C42', 'B42', 'C67', 'B36', 'C10'\}} \\
\texttt{ARGUMENT: \{'I am pleased to accept this revised proposal, as I believe it strikes an excellent balance between effort and importance. The inclusion of C10 adds significant importance without exceeding the effort limit, and I think this proposal is a very good compromise. I am satisfied with the outcome of the negotiation and appreciate the other player\textbackslash{}'s willingness to revise and improve the proposal.'\}} \\
            }
        }
    }
     \\ \\

    \theutterance \stepcounter{utterance}  
    & & & \multicolumn{2}{p{0.3\linewidth}}{
        \cellcolor[rgb]{0.9,0.9,0.9}{
            \makecell[{{p{\linewidth}}}]{
                \texttt{\tiny{[GM$|$GM]}}
                \texttt{game successful} \\
            }
        }
    }
    & & \\ \\

    \theutterance \stepcounter{utterance}  
    & & & \multicolumn{2}{p{0.3\linewidth}}{
        \cellcolor[rgb]{0.9,0.9,0.9}{
            \makecell[{{p{\linewidth}}}]{
                \texttt{\tiny{[GM$|$GM]}}
                \texttt{end game} \\
            }
        }
    }
    & & \\ \\

\end{supertabular}
}

\end{document}
